\newpage
\thispagestyle{empty}
\phantom{.}
\newpage
\thispagestyle{empty}
\begin{center}
    {{\usefont{OT1}{cmr}{b}{n}\LARGE\selectfont{Anotační list}}\par}
    \vspace*{6mm}
    \begin{tabular}{@{}ll@{}}
        \textbf{Jméno autora:} & \textbf{Štěpán Borovec} \\
        \textbf{Název práce česky:} & \textbf{Dekarbonizace v rafinérském průmyslu} \\
        \textbf{Název práce anglicky:} & \textbf{Decarbonisation in the oil refinery} \\
        \textbf{Jazyk práce:} & čeština \\
        \textbf{Akademický rok:} & 2023/2024 \\
        \textbf{Studijní program:} & Teoretický základ strojního inženýrství \\
        \textbf{Ústav:} & Ústav procesní a zpracovatelské techniky \\
        \textbf{Vedoucí práce:} & doc. Ing. Radek Šulc, Ph.D. \\
        \textbf{Rozsah práce:} & \ref{TotPages} stran\\
         & \totalfigures\ obrázků\\
         & \totaltables\ tabulek\\
        \textbf{Klíčová slova:} & vodík, dekarbonizace, emise, rafinérie\\
        \textbf{Key words:} & hydrogen, decarbonization, emissions, refinery\\
    \end{tabular}
\end{center}
\vspace*{3mm}
\paragraph*{\textbf{Abstrakt:}}
Cílem této bakalářské práce je přispět k lepšímu pochopení potenciálu vodíku jako nosiče energie a jeho využití v rafinériích. Bakalářská práce je zaměřena na analýzu aktuálních projektů dekarbonizace rafinérií pomocí zeleného vodíku a návrh vlastního výpočtu pro zlepšení udržitelnosti rafinérií. V první kapitole jsou popsány základní charakteristiky vodíku, jeho fyzikální a chemické vlastnosti a různé typy vodíku v závislosti na způsobu výroby. Následující kapitoly se věnují konkrétním aplikacím vodíku v průmyslu, včetně jeho využití v chemickém, metalurgickém a dopravním sektoru, a podrobně jsou popsány procesy výroby vodíku v rafinériích. Závěrečná část práce je zaměřena na aktuální projekty dekarbonizace a možnosti využití zeleného vodíku, včetně návrhu vlastní studie na toto téma.
\vspace*{3mm}
\paragraph*{\textbf{Abstract:}}
The aim of this bachelor thesis is to contribute to a better understanding of the potential of hydrogen as an energy carrier and its implementation in refineries. The bachelor thesis focuses on the analysis of current projects of decarbonization of refineries using green hydrogen and the design of a custom calculation to improve the sustainability of the refinery. The first chapter describes the basic characteristics of hydrogen, its physical and chemical properties and the different types of hydrogen depending on the production method. Subsequent chapters discuss specific applications of hydrogen in industry, including its use in the chemical, metallurgical and transport sectors, and detail the hydrogen production processes in refineries. The final part of the thesis focuses on current decarbonisation projects and the potential use of green hydrogen, including the design of a self-study on this topic.